\documentclass{beamer}


%\beamertemplateshadingbackground{yellow!100}{white}
\usepackage[english]{babel}
\usepackage{amsfonts,amsmath}
\usepackage{graphicx}
\usepackage{sansmathaccent}
\pdfmapfile{+sansmathaccent.map}

%\usetheme{Warsaw}
\usetheme[secheader]{Madrid}

\usepackage{multimedia}
\logo{\includegraphics[height=0.5cm]{imi.jpg}}

\newcommand{\be}{\begin{equation}}
\newcommand{\ee}{\end{equation}}
\newcommand{\rf}[1]{(\ref{#1})}
\newcommand{\RR}{\mathbb{R}}
\newtheorem{thm}{Theorem}
\newtheorem{lm}{Lemma}

%\def\ra#1\{\renewcommand{\arraystretch}{#1}
\begin{document}
 %[propagating wave solutions to the 2D BPE]
\title{Comparison between two numerical methods for solution of 2D Boussinesq paradigm equation}

\author[K. Angelow]{{\underline{Krassimir Angelow}}, Natalia Kolkovska}
\institute[IMI -- BAS]{Institute of Mathematics and Informatics\\ Bulgarian Academy of Sciences, Sofia, Bulgaria,\\ e-mail: angelow@math.bas.bg}
\date[2021]{AMITANS, June 20-25  2021,  Albena, Bulgaria}


%---------- frame 01 ----------------
\begin{frame}
 \titlepage
 \begin{center}
  Supported by Bulgarian National Science Fund Grant K$\Pi$-06-H22/2
  \end{center}
\end{frame}

%---------- frame 02 ----------------
\begin{frame}
\tableofcontents 
\setbeamertemplate{table of contents shaded}[default]
\section{Boussinesq Paradigm Equation}
\section{Quadrature Formulas}
\section{Numerical Methods}
\subsection{Taylor Series (TS) Approach}
\subsection{Energy Saving Method with Conservative Scheme}

%---------- frame 06 ----------------
\section{Convergence}

%---------- frame 07 ----------------
\section{Comparison Between TS and Conservative Scheme}
%\tableofcontents 
\end{frame}

%---------- frame 03 ----------------
\begin{frame}
\frametitle{Boussinesq Paradigm Equation}


We study the problem

\be\label{problem}
\beta(E-\Delta) \frac{\partial^2 u}{\partial t^2}=
 \beta \Delta u -\Delta^2 u -\alpha \beta \Delta (u^2)
\ee
where $\alpha>0$, $\beta>0$  are dispersion parameters, and $\Delta$ is the Laplace operator.
\end{frame}



%---------- frame 04 ----------------
\begin{frame}
\frametitle{Quadrature Formulas for the Integral and Energy}
\framesubtitle{Integral and Energy}
The integral 

\begin{equation}\label{int}
D(u)=\int_{R^2} u(x,y)dx dy
\end{equation}

and energy:
\begin{align}\label{ex-en}
E(u):=&\int_{R^2} u_t \left((A^{-1}+E)u_t\right) dxdy+
\beta \int_{R^2} u^2 dxdy \nonumber\\
+& \int_{R^2}u \left(A u\right) dxdy
-\frac{2 \alpha \beta}{3} \int_{R^2} u^3 dxdy =const
\end{align}
of the continuous problem \rf{problem}. Here $E(u)$ is the exact energy of problem \rf{problem} and $Au=-\Delta u$.
\end{frame}

%---------- frame 05 ----------------
\begin{frame}
\frametitle{Quadrature Formulas for the Integral and Energy}
\framesubtitle{2D trapezoidal formula (with global $O(h_1^2+h_2^2)$ error)}

Let's assign the X,Y dimensions of the computational domain with $N_1$ and $N_2$.
Then the approximation of the integral \eqref{int} with global $O(h_1^2+h_2^2)$ error is

\begin{align}\label{quadr2}
D_h(u_{i,j}) =& \sum_{i=1}^{N_1-1} \sum_{j=1}^{N_2-1} h_1 h_2 u_{i,j}
+\frac{h_1}{2}\sum_{i=0} \sum_{j=1}^{N_2-1} h_2 u_{i,j}
+\frac{h_1}{2}\sum_{i=N_1} \sum_{j=1}^{N_2-1} h_2 u_{i,j} \nonumber\\
+&\frac{h_2}{2}\sum_{j=0} \sum_{i=1}^{N_1-1} h_1 u_{i,j}
+\frac{h_2}{2}\sum_{j=N_2} \sum_{i=1}^{N_1-1} h_1 u_{i,j}
\nonumber\\
+&\frac{1}{4}h_1 h_2 \left(u_{0,0}+u_{N_1,0}+u_{N_1,N_2}+u_{0,N_2}
\right).
\end{align}

\end{frame}

\begin{frame}
\frametitle{Quadrature Formulas for the Integral and Energy}
\framesubtitle{2D Simpson's rule (with global $O(h_1^4+h_2^4)$ error)}
Assume that $N_1=2k$, $N_2=2 l$.
For every $m=0,1,2,\cdots N_2$ we compute 

$$D_m= \frac{h_1 }{3} 
\left\{ u_{0,m}+u_{N_1,m}+ 4 \sum_{i=1}^{\frac{N_1}{2}}   u_{2i-1,m}
 +2 \sum_{i=1}^{\frac{N_1}{2}-1} u_{2i,m} \right\}$$


Then 

\begin{equation}\label{quadr4}
D_h(u)=\frac{h_2 }{3} 
\left\{ D_{0}+D_{N_2}+ 4 \sum_{j=1}^{\frac{N_2}{2}}   D_{2j-1}
 +2 \sum_{j=1}^{{\frac{N_2}{2}}-1} D_{2j} \right\}
\end{equation}
is the approximation of the integral \eqref{int} with global $O(h_1^4+h_2^4)$ error

\end{frame}

\begin{frame}
\frametitle{Quadrature Formulas for the Integral and Energy}
\framesubtitle{2D Boole's rule (with global $O(h_1^6+h_2^6)$ error)}

Assume that $N_1=4k$, $N_2=4 l$.
For every $m=0,1,2,\cdots N_2$ we compute

\begin{align*}
D_m =& \frac{2h_1}{45} 
\left\{
7u_{0,m}+7u_{N_1,m}+32 \sum_{i=1}^{\frac{N_1}{2}}u_{2i-1,m}
+12\sum_{i=1}^{\frac{N_1}{4}}u_{4i-2,m}
+14 \sum_{i=1}^{\frac{N_1}{4}-1}u_{4i,m}
\right\}
\end{align*}

then \eqref{quadr6-2D} is the approximation to the integral \eqref{int} with global $O(h_1^6+h_2^6)$ error

\begin{align}\label{quadr6-2D}
&D_h(u) =
\frac{2h_2}{45} 
\left\{
7D_{0}+7D_{N_2}+32 \sum_{j=1}^{\frac{N_2}{2}}D_{2j-1}
+12\sum_{j=1}^{\frac{N_2}{4}}D_{4j-2}
+14 \sum_{j=1}^{\frac{N_2}{4}-1}D_{4j}
\right\}
\end{align}
\end{frame}

\begin{frame}
\frametitle{Numerical Methods}
\framesubtitle{Taylor Series}

Suppose we have functions $y$ and $f$ with the following properties
\begin{equation}
y:I \rightarrow \RR, \quad f:I \times \Omega \rightarrow \RR \nonumber
\end{equation}
where $y \in C^1(I,\RR)$ and $f$ is sufficiently smooth on $I \times \Omega$. Then, $f$ represent the phase space and $y$ is the unknown function to the differential equaiton:
\begin{equation}
\dot{y}(t) = f(t, y(t)).
\end{equation}
By repeated differentiation we can find each function

\begin{equation}
\frac{d^s}{dt^s}y(t) = \frac{d^{s-1}}{dt^{s-1}}f(t, y(t))
\end{equation}

and evaluate it at $t=t_0$ for each derivative.
\end{frame}

\begin{frame}
\frametitle{Numerical Methods}
\framesubtitle{Taylor Series}

The order $s$ formula for computing $y(t_{n+1})=$ $y(t_{n}+h)$ using these functions, evaluated at $t=t_{n}$ and $y=y_{n}$, is
\begin{eqnarray}
y_{n+1} = y_{n} + h f(t_{n}, y_{n}) + \frac{h^2}{2!} \frac{d^2}{dt^2} f(t_{n}, y_{n}) + ... + \frac{h^s}{s!} \frac{d^s}{dt^s} f(t_{n}, y_{n}), &\nonumber\\ 
f_i(t_l, y_l), y_l(t_l), t_l \in \RR, 1 \leq l \leq n, 1 \leq i \leq s
\end{eqnarray}
where we already have the pair $(t_0, y_0)$ from the initial condition of the differential equation.

\end{frame}

\begin{frame}
\frametitle{Numerical Methods}
\framesubtitle{Energy Saving Method with Conservative Scheme}

The approximation of the differential operators is defined as:
\begin{equation}
\frac{\partial^2 u}{\partial t^2}(x_i, y_j, t_k ) = \frac{ u^{(k+1)}_{i, j} - 2u^{(k)}_{i,j} + u^{(k-1)}_{i,j} }{\tau^2} + O(\tau^2) 
\end{equation}

\begin{equation}
\frac{\partial^2 u}{\partial x^2}(x_i, y_j, t_k ) = \frac{ u^{(k)}_{i+1, j} - 2u^{(k)}_{i,j} + u^{(k)}_{i-1,j} }{h_1^2} + O(h_1^2) 
\end{equation}

\begin{equation}
\frac{\partial^2 u}{\partial y^2}(x_i, y_j, t_k ) = \frac{ u^{(k)}_{i, j+1} - 2u^{(k)}_{i,j} + u^{(k)}_{i,j-1} }{h_2^2} + O(h_2^2) 
\end{equation}


\begin{equation}
\Delta_h = \frac{ u^{(k)}_{i+1, j} - 2u^{(k)}_{i,j} + u^{(k)}_{i-1,j} }{h_1^2} + \frac{ u^{(k)}_{i, j+1} - 2u^{(k)}_{i,j} + u^{(k)}_{i,j-1} }{h_2^2}
\end{equation}

\end{frame}

\begin{frame}
\frametitle{Numerical Methods}
\framesubtitle{Energy Saving Method with Conservative Scheme}
The grid function is defined as:
\begin{equation}
(I-\Delta_h)\frac{ u^{(k+1)}_{i, j} - 2u^{(k)}_{i,j} + u^{(k-1)}_{i,j} }{\tau^2} = (\Delta_h - \Delta_h^2)u^{(k)}_{i,j} + \Delta_h(g(u^{(k)}_{i,j})).
\end{equation}
%
where the non-linear term $g$ is defined as:
\begin{align}
g(u^{(k)}_{i,j})=&\frac{\alpha \beta} { 3 } \left( (u^{(k+1)}_{i,j})^2 + (u^{(k-1)}_{i,j})(u^{(k+1)}_{i,j}) + (u^{(k-1)}_{i,j})^2 \right) + \nonumber\\
+&\frac{ (\beta - 1 )}{ 2 }\left( u^{(k+1)}_{i,j} + u^{(k-1)}_{i,j} \right).
\end{align}


\end{frame}


\begin{frame}
\frametitle{XXXXXXXXXXXX}

\end{frame}


\end{document}

